% purpose: 读者应能回答"控制体如何按环形分割、界面通量写在哪里、为何中心零通量自动成立、表面未知值由哪两个通量消元得到"
% topology: geometry (轴对称环形控制体) + derivation (通量消元)
% reading_direction: coordinate-based（径向由中心向外）+ 局部放大
% node_roles: 环形控制体、界面、格心、表面、外侧空气
% edge_roles: 界面扩散通量、内侧离散通量、外侧对流通量
% visual_encoding: 实线=界面与实体边界；斜线阴影=当前控制体；虚线箭头=质流；双箭头尺寸线=几何距离；无颜色
% style_family: C1（纯黑白线稿，与全篇 mono 流程图一致）
% annotation_policy: minimal   result_policy: symbolic   final_width: text   print_mode: grayscale
% layout_rule: 扇形内只放短标签，含公式的说明块一律放到扇形外；右侧说明块共用同一左缘 \Xc
\documentclass[border=2pt]{standalone}
\usepackage{ctex}
\usepackage{amsmath}
\usepackage{tikz}
\usetikzlibrary{arrows.meta,calc,patterns,decorations.pathmorphing}
\begin{document}
\begin{tikzpicture}[
  font=\normalsize,
  solidline/.style={line width=0.9pt},
  facelin/.style={line width=0.6pt},
  thinline/.style={line width=0.5pt},
  massflow/.style={-{Stealth[length=4.2pt,width=3.0pt]},line width=0.7pt,densely dashed},
  lead/.style={line width=0.4pt,densely dotted},
  dimline/.style={{Stealth[length=3.4pt,width=2.6pt]}-{Stealth[length=3.4pt,width=2.6pt]},line width=0.5pt}
]
% ---- 主图：四分之一截面。环形控制体的"图上面积"即体积权，中心退化为零面积。
\def\Rs{4.20}          % eta=1 对应的绘图半径
\def\fA{0.75}\def\fBi{2.55}\def\fBm{3.10}\def\fBo{3.65}
\def\Xc{4.80}          % 右侧说明块统一左缘：全部 anchor=north west，避免左边界参差

% 当前控制体 i：环形扇形阴影（先画底层，后画界面线，避免遮盖）
\fill[pattern=north east lines]
      (\fBi,0) -- (\fBm,0) arc[start angle=0,end angle=90,radius=\fBm]
   -- (0,\fBi) arc[start angle=90,end angle=0,radius=\fBi] -- cycle;

% 界面弧（细实线）与两条对称半径
\foreach \r in {\fA,\fBi,\fBm,\fBo}{
  \draw[facelin] (\r,0) arc[start angle=0,end angle=90,radius=\r];
}
\draw[thinline] (0,0) -- (\Rs,0);
\draw[thinline] (0,0) -- (0,\Rs);
% 外表面（粗实线=给定物理边界）
\draw[solidline] (\Rs,0) arc[start angle=0,end angle=90,radius=\Rs];

% ---- 格心（落在水平半径上）与省略的中间控制体
\foreach \x in {0.375,2.825,3.375,3.925}{ \fill (\x,0) circle (1.5pt); }
\node[inner sep=1pt] at (1.65,0.30) {$\cdots$};

% 格心记号排在水平半径下方，用点引线相连，避免压住阴影与界面弧
\foreach \x/\lb in {0.375/{$C_1$},2.825/{$C_i$},3.925/{$C_N$}}{
  \draw[lead] (\x,-0.06) -- (\x,-0.42);
  \node[anchor=north,inner sep=1.5pt] at (\x,-0.42) {\lb};
}
\draw[lead] (3.375,-0.06) -- (3.375,-0.98);
\node[anchor=north,inner sep=1.5pt] at (3.375,-0.98) {$C_{i+1}$};

% ---- 界面与端点坐标记号（贴在竖直半径外侧的空白区，互不相压）
\node[anchor=east,inner sep=2pt] at (-0.06,\fBi) {$\eta_{i-1/2}$};
\node[anchor=east,inner sep=2pt] at (-0.06,\fBm) {$\eta_{i+1/2}$};
\node[anchor=north east,inner sep=2pt] at (-0.02,-0.02) {$\eta=0$};

% ---- 受热/传质表面：标签移到左上空白角，用引线指到外弧，兼顾填白与不压线
\node[anchor=south west,align=left,inner sep=2pt] at (-0.90,4.44) {受热/传质表面 $\eta=1$};
\draw[lead] (0.62,4.40) -- (1.98,3.72);

% ---- 阴影环形即控制体 i：扇形内只放不含公式的短标签（宽度小于无弧环带）
\node[inner sep=2pt] at (1.15,1.30) {控制体 $i$};
\draw[lead] (1.58,1.46) -- (2.06,2.06);

% ---- 界面扩散通量（虚线箭头沿径向跨过界面弧）
% 加长并加粗：原 11pt 长、0.7pt 的虚线压在 0.6pt 斜线阴影上会被读成纹理的一部分。
% 仍保持"虚线=质流"的编码不变，只把线重提到明显高于阴影线，并保留两箭头之间的断口。
\draw[massflow,line width=0.9pt] (2.11,0.76) -- (2.58,0.94);
\draw[massflow,line width=0.9pt] (2.73,0.99) -- (3.24,1.185);

% ---- 局部放大：最外控制体与表面（局部曲率可忽略，故画成平直薄层）
% 放大引线由主图表面弧的两点拉向放大条的左端，指明放大的是哪一段
\draw[lead] (3.95,1.44) -- (5.20,3.45);
\draw[lead] (3.44,2.41) -- (5.20,4.55);
% 坐标写成绝对值而不用 shift scope，便于自动检查真实读到节点位置
\draw[facelin] (5.20,3.45) -- (5.20,4.55);
\draw[thinline] (5.20,3.45) -- (8.20,3.45);
\draw[thinline] (5.20,4.55) -- (8.20,4.55);
\draw[solidline] (8.20,3.45) -- (8.20,4.55);
% 放大条标题排在"内侧离散扩散"之上一行，两者同为 anchor=south，必须留出整行间距
\node[anchor=south west,inner sep=2pt] at (5.20,5.26) {表面附近放大};
% 格心与待求表面值：格心不落在表面上，这是本题的分水岭
\fill (6.20,4.00) circle (1.5pt);
\node[anchor=north,inner sep=2pt] at (6.20,3.92) {$C_N$};
\draw[thinline] (8.20,4.00) circle (1.9pt);
\node[anchor=south,inner sep=2pt] at (8.35,5.30) {待求 $C_s$};
\draw[lead] (8.32,5.26) -- (8.22,4.08);
% 几何距离：格心到表面只有半个网格（尺寸线画在薄层之外）
\draw[dimline] (6.20,3.20) -- (8.20,3.20);
\node[anchor=north,inner sep=2pt] at (7.20,3.18) {$R\Delta\eta/2$};
% 两侧通量在表面相遇
\draw[massflow] (6.30,4.30) -- (8.10,4.30);
\draw[massflow] (8.30,4.30) -- (9.05,4.30);
\node[anchor=south,inner sep=2pt] at (6.28,4.62) {内侧离散扩散};
\node[anchor=north west,align=left,inner sep=2pt] at (8.35,4.16) {外侧\\对流传质};

% ---- 右侧说明列：三块共用左缘 \Xc，各自贴近它所解释的对象
% 表面值反解（紧贴放大图下方）
\node[anchor=north west,align=left,inner sep=2pt] at (\Xc,2.62)
     {表面两侧通量相等，消元得\\[2pt]
      $C_s=(\beta C_N+h_mC_a)/(\beta+h_m)$\\[2pt]
      其中 $\beta=2D_N/(R\Delta\eta)$};
% 界面通量定义（与主图中跨界面的虚线箭头等高）
\node[anchor=north west,align=left,inner sep=2pt] at (\Xc,0.55)
     {跨界面弧的扩散通量\\[2pt]
      $J_{i+1/2}=\eta_{i+1/2}D_{i+1/2}\dfrac{C_{i+1}-C_i}{\Delta\eta}$\\[2pt]
      界面物性取调和平均};
% 守恒平衡：控制体的变化率只由两个界面通量之差给出
% 守恒平衡下移，填住右下空白带，同时与上一块留出整行间距
\node[anchor=north west,align=left,inner sep=2pt] at (\Xc,-1.70)
     {控制体 $i$ 的守恒平衡\\[2pt]
      $R_0^{2}V_i\dfrac{\mathrm{d}C_i}{\mathrm{d}t}=J_{i+1/2}-J_{i-1/2}$};

% ---- 底部左侧：体积权与中心退化。含公式，故放在扇形之外的空白带
\node[anchor=north west,align=left,inner sep=2pt] at (-0.95,-1.62)
     {画面环形面积 $\propto$ 体积权\\[1pt]
      $V_i=(\eta_{i+1/2}^{2}-\eta_{i-1/2}^{2})/2$\\[2pt]
      中心面积权 $\eta_{-1/2}=0$，\\[1pt]
      轴心零通量自动满足};
\end{tikzpicture}
\end{document}
